
% ================================================================
% CHAPTER 1: INTRODUCTION
% ================================================================

\section{System Overview}

\subsection{General System}

Imagine asking an AI assistant: ``Who is the current Prime Minister of India?''
The AI confidently replies ``Narendra Modi.'' Your document collection, however,
contains a more recent article saying a different leader has taken over.
Which source should the system trust? That tension between what the model
``remembers'' from training and what the documents actually say is the
core problem TruthfulRAG v5 was built to solve.

TruthfulRAG v5 is a Knowledge-Graph based question-answering pipeline that builds
a structured map of every fact in the document collection, spots contradictions,
decides which side of the conflict is more trustworthy (based on age and source
agreement), and hands the LLM a clean, conflict-free prompt. A confidence
percentage accompanies every answer. Everything runs locally — no cloud, no
subscription, using Qwen2.5-7B-Instruct via Ollama and Neo4j Community Edition.

\subsection{Overall Project Flow}

The pipeline runs in four sequential stages:
\begin{enumerate}\itemsep0pt\parsep0pt
    \item \textbf{Schema Inference}: The LLM samples a handful of documents and
    identifies entity and relation types automatically — no developer configuration.
    \item \textbf{Graph Construction (Module A)}: Facts are extracted and stored
    as typed, timestamped, weighted edges in Neo4j.
    \item \textbf{Graph Retrieval (Module B)}: Personalised PageRank finds
    structurally nearby facts; temporal decay penalises older information.
    \item \textbf{Conflict Resolution (Module C)}: Contradicting facts are
    compared by score gap and entropy; the weaker is discarded with a logged reason.
\end{enumerate}

\subsection{Key Features}

Five properties distinguish this system from standard RAG:
\begin{enumerate}\itemsep0pt\parsep0pt
    \item \textbf{Zero configuration}: Schema is discovered automatically from the corpus.
    \item \textbf{Time-awareness}: Temporal decay $e^{-0.08 \times \text{age}}$ is
    applied at retrieval so newer facts outrank older ones on the same topic.
    \item \textbf{Calibrated honesty}: A three-factor confidence score replaces
    the implicit 100\,\% certainty of standard RAG.
    \item \textbf{Explainability}: Every conflict removal is logged with the
    numerical reason.
    \item \textbf{Domain independence}: The same code handles medical, legal,
    space-science, and political corpora without modification.
\end{enumerate}

% ----------------------------------------------------------------
\section{Problem Statement}

\subsection{Rationale for Research}

Large language models have a frozen memory. Everything they know was baked in
during training, and if the world changed after the cutoff date, the model does
not know, but it will often answer confidently regardless. Standard RAG was
meant to close that gap by supplying documents at query time, but real document
collections are messy: a hospital's records may contain both a 1990 guideline
and a 2023 update for the same drug. Standard RAG passes both to the LLM,
which has no principled way to choose and either hedges or picks arbitrarily.
The motivating observation is simple: a careful human reader would compare
dates, weight the newer source, and check for corroboration. A software system
can and should do the same.

\subsection{Gaps in Existing Tools}

Four concrete gaps exist in current open-source RAG tools:
\begin{enumerate}\itemsep0pt\parsep0pt
    \item No tool detects conflicts between retrieved documents — LangChain,
    LlamaIndex, and GraphRAG pass whatever they retrieve directly to the LLM.
    \item KG-RAG systems require manual entity/relation specification; switching
    domains means rewriting configuration files.
    \item No system attaches a confidence score — every response is presented
    with equal certainty regardless of evidence quality.
    \item Conflict resolution, when it occurs at all, is silent — no audit trail
    explains what was kept, what was discarded, and why.
\end{enumerate}

\subsection{Challenges in Document-Based QA}

Three properties of real-world document collections consistently challenge
retrieval pipelines:
\begin{enumerate}\itemsep0pt\parsep0pt
    \item \textbf{Facts are time-sensitive.} A system with no concept of time
    answers current-state questions using whatever it retrieves first, regardless of age.
    \item \textbf{Agreement across sources matters.} Seven independent documents
    confirming a fact constitute far stronger evidence than one; retrieval
    scores should reflect this.
    \item \textbf{Domains vary wildly.} Any manual setup step is a barrier that
    limits real-world deployment across heterogeneous document collections.
\end{enumerate}

\subsection{Project Objectives}

The five concrete objectives of this project were:
\begin{enumerate}\itemsep0pt\parsep0pt
    \item Faithfully re-implement the KG-RAG v4 pipeline (arXiv:2511.10375)
    in Python using LangChain, Ollama, and Neo4j.
    \item Extend that baseline with nine novel contributions addressing the
    gaps above, particularly temporal scoring and corroboration weighting.
    \item Eliminate the last remaining manual step in v4 (schema configuration)
    by adding automatic domain inference before graph construction.
    \item Validate the system on four real-world domains: Indian political
    history, Indian science, Indian law, and space science.
    \item Evaluate performance against both the v4 baseline and standard
    LangChain RAG using benchmark datasets and conflict-detection metrics.
\end{enumerate}

% ----------------------------------------------------------------
\section{Motivation}

\subsection{Why This Matters}

A system that says ``I am 42\,\% confident in this answer'' is more useful in
a hospital or a courtroom than one that always sounds completely sure.
Overconfident AI answers in high-stakes settings are not just wrong, they are
dangerous. This project sets out to show that honest, conflict-aware QA does
not require a well-funded research lab: it can be built by a final-year
undergraduate student on a laptop using freely downloadable, open-source tools,
without sending a single document to an external server. The organisations
that most need trustworthy AI — public hospitals, lower courts, government
archives — are exactly those that cannot afford cloud subscriptions or
proprietary knowledge-base software.



